% Part: many-valued-logic
% Chapter: infinite-valued-logics
% Section: introduction

\documentclass[../../../include/open-logic-section]{subfiles}

\begin{document}

\olfileid[hi]{mvl}{inf}{int}

\olsection{परिचय}

किसी आव्यूह के सत्य-मूल्यों की संख्या परिमित होना आवश्यक नहीं है। अनंत रूप
से अनेक सत्य-मूल्यों के समुच्चय के लिए एक स्वाभाविक विकल्प $0$ और $1$ के
बीच की परिमेय संख्याओं का समुच्चय $V_\infty = [0,1] \cap \Rat$ है, अर्थात्
\begin{align*}
    V_\infty & = \Setabs{\frac{n}{m}}{n,m \in \Nat \text{ और } n\le m}.
\intertext{इस अनंत सत्य-मूल्य समुच्चय पर विचार करते समय निम्न उपसमुच्चयों पर
भी विचार करना प्रायः उपयोगी होता है:}
V_m & = \Setabs{\frac{n}{m-1}}{n \in \Nat \text{ और } n\le m}.
\intertext{उदाहरणार्थ, $V_5$ समान अंतराल पर स्थित $5$ सत्य-मूल्यों का
समुच्चय है:}
V_5 & = \{0, \frac{1}{4}, \frac{1}{2}, \frac{3}{4}, 1\}.
\end{align*}
इन सत्य-मूल्य समुच्चयों पर आधारित तर्कशास्त्रों में सामान्यतः केवल $1$
निर्दिष्ट होता है, अर्थात् $V^+ = \{1\}$। दूसरे शब्दों में, हम $1$ को
(निरपेक्ष) सत्यता और $0$ को निरपेक्ष असत्यता की भूमिका देते हैं, किंतु
!!{formula}ें $V$ का कोई भी मध्यवर्ती मूल्य ले सकती हैं।

इसके अतिरिक्त, $0$ और $1$ के बीच की सभी \emph{वास्तविक} संख्याओं के
समुच्चय $V_{[0,1]} = [0,1]$, या $[0,1]$ के अन्य अनंत उपसमुच्चयों पर भी
विचार किया जा सकता है। इस सत्य-मूल्य समुच्चय वाले तर्कशास्त्र प्रायः
\emph{फ़ज़ी} कहलाते हैं।

\end{document}
